\documentclass[11pt, a4paper]{article}
\input{preamble.tex}
\switchlang{ru}
\begin{document}

\title{1989 -- Numonics Manager Mouse Cordless, 2nd generation}
\date{}
\maketitle
\selectlanguage{russian}

Мышь Numonics Manager Mouse была второго поколения является переработанным вариантом исходной Manager Mouse, выпущенной в 1985 году компанией Torrington, крупным международным производителем подшипников со штаб-квартирой в Коннектикуте \cite{torringtonAdd}. Поскольку для Torrington выпуск мышей был достаточно далек от их основной сферы деятельности, они охотно лицензировали Manager Mouse другим компаниям, например Maynard Electronics \cite{admaynard} и Numonics Corporation \cite{adnumonics}. Numonics специализировалась на производстве прецизионных измерительных приборов, а также получила известность как производитель компьютерной периферии для САПР, такой как дигитайзеры и плоттеры \cite{numonics}. Сначала она включила мышь, разработанную Torrington, в спектр устройств, доступных под своим брэндом, а затем в 1989 году представила на рынке обновленную модель мыши под тем же названием "--- Numonics Manager Mouse. Обновленный вариант не встречается под брэндами других производителей, будучи эксклюзивным устройством и предназначался для САПР и работы с графикой \cite{buxtonNumonicsCardboard}, что логично, учитывая направление деятельности компании. Код FCC ID, присутствующий на корпусе, также принадлежит Numonics Corporation, но тем не менее участие Torrington в качестве фактического производителя устройства также вполне возможно.

\begin{figure}[h]
   \centering
    \includegraphics[scale=0.9]{1989_numonics_manager_mouse_cordless_2gen/pic_30.jpg}
    \caption{Manager Mouse}
    \label{fig:NumonicsManagerMouse2Pic}
\end{figure}

Как и Manager Mouse первого поколения, эта мышь также появилась на рынке в двух вариантах: с проводным (Manager Mouse 3S) и с беспроводным инфракрасным подключением (Manager Mouse Cordless) \cite{italianReview}. Модель была доступна для IBM-совместимых компьютеров и для Apple Macintosh (последняя "--- только в проводном варианте) \cite{buxtonNumonicsCardboard}. Вслед за первой моделью Manager Mouse, данная мышь сохранила и необычный способ регистрации движения на основе колес.

Датировку мыши можно определить по коду FCC ID, который, согласно базе данных Федеральной комиссии по связи США, показывает, что мышь была выпущена на рынок в 1989 году, о чем свидетельствует и почтовый штамп на отсканированном варианте упаковки \cite{buxtonNumonicsCardboard}.

\section*{Внешний вид}

Корпус мыши выполнен из белого пластика. Как и у мыши первого поколения, в основание корпуса Manager Mouse 2-го поколения вписан прямоугольник "--- однако на этом сходство заканчивается. Верхняя часть корпуса имеет обтекаемую <<аэродинамическую>> форму, ее продольное сечение представляет собой асимметричный колоколообразный профиль с подъёмом в области запястья пользователя. Три крупные кнопки одного цвета с корпусом полностью интегрированы в его форму и не имеют каких-либо выступов (средняя кнопка снабжена круглой выпуклостью для облегчения тактильной идентификации). Какие-либо другие элементы на верхней стороне корпуса отсутствуют, за исключением названия <<Manager Mouse>>, нанесенного краской возле левой кнопки. Передняя (дальняя от пользователя) часть корпуса содержит инфракрасный излучатель: поэтому она закрыта выпуклым <<бампером>> из глянцевого черного пластика (у проводного варианта мыши корпус полностью белый), который в ходе информационном обмене подсвечивается светодиодами, приобретая полупрозрачный голубоватый оттенок \cite{italianReview}. В центре бампера (там, где у проводного варианта расположен выход провода мыши) находится отверстие, через которое в беспроводную мышь вставляется штекер зарядного кабеля.

\begin{figure}[h]
    \centering
    \includegraphics[scale=0.9]{1989_numonics_manager_mouse_cordless_2gen/top_15.jpg}
    \includegraphics[scale=0.9]{1989_numonics_manager_mouse_cordless_2gen/bottom_15.jpg}
    \caption{Manager Mouse, вид сверху и снизу}
    \label{fig:NumonicsManagerMouse2TopAndBottom}
\end{figure}

Нижняя часть мыши показывает большое сходство с Torrington Manager Mouse (рис. \ref{fig:NumonicsManagerMouse2TopAndBottom}): на ней можно видеть наклейку с информационными данными, круглые накладки из низкофрикционного материала, облегчающие скольжение мыши, а также два маленьких колеса в форме усечённого конуса, предназначенные для регистрации движения вместо более традиционного для механических мышей шара. Это решение восходит к первым мышам Дугласа Энгельбарта: колеса расположены ортогонально друг другу, чтобы при продольном и поперечном движении одно из них вращалось, а другое проскальзывало. Но, по сравнению с мышью Энгельбарта, колеса Manager Mouse существенно уменьшены в диаметре, а их оси расположены под углом приблизительно 15 градусов к горизонтальной плоскости \cite{italianReview}, что должно обеспечивать регистрацию диагональных перемещений.

\section*{Размеры и эргономика}

\begin{figure}[h]
    \centering
    \includegraphics[scale=0.51]{1989_numonics_manager_mouse_cordless_2gen/size_15.jpg}
    \caption{Manager Mouse на размерном коврике с шагом сетки 1~см}
    \label{fig:NumonicsManagerMouse2Size}
\end{figure}

Корпус имеет размер, типичный для второй половины 1980-х годов (рис. \ref{fig:NumonicsManagerMouse2Pic}, \ref{fig:NumonicsManagerMouse2Size}).
Однако благодаря тому, что по сравнению с первым поколением Manager Mouse его геометрия претерпела значительные изменения, форму данной мыши можно назвать одной из наиболее эргономичных среди мышей конца 80-х годов. Это касается и более широких кнопок, вписанных в форму корпуса, и выпуклой опоры для ладони. К тому же корпус мыши является симметричным, одинаково удобным как для левшей, так и для правшей.

В обзоре \cite{italianReview} отмечается, что особенности механизма регистрации движения обеспечивают весьма точное позиционирование указателя, однако конструкция колес все же имеет недостатки: они не идеально реагируют на диагональные движения (таким образом, наклонное положение осей, призванное решать эту проблему, справляется с задачей не идеально).

\begin{figure}[h]
    \centering
    \includegraphics[scale=0.51]{1989_numonics_manager_mouse_cordless_2gen/hand_15.jpg}
    \caption{Manager Mouse с моделью руки человека}
    \label{fig:NumonicsManagerMouse2Hand}
\end{figure}

\section*{Подключение}

И Torringon/Numonics в 1986 году, и даже Numonics в 1989 году рекламировали беспроводной вариант Manager Mouse как <<первую в мире беспроводную мышь>> (в действительности первой мышью таким интерфейсом является инфракрасная мышь, разработанная Logitech в 1984 году для рабочих станций Metaphor Computer Systems). Как бы там ни было, как и проводной вариант Manager Mouse, инфракрасный приемник беспроводной версии подключается к компьютеру через порт RS-232 (поддерживаются протоколы Microsoft и Mouse Systems). Кабель приемника оснащен 9-контактным разъемом последовательного интерфейса, в который также подключается 9-вольтный блок питания. Кроме того в комплект входят переходник на 25-контактный разъём RS-232 и зарядный кабель. Для зарядки мыши этим кабелем необходимо соединить разъем на приемнике и гнездо в передней части мыши. В обзоре \cite{italianReview} подчеркивается, что для зарядки не требуется соединение с последовательным портом, а также что время полной зарядки составляет не менее 15 часов, после чего устройство может работать около 8 часов. При этом оно автоматически переходит в режим ожидания через одну минуту бездействия, после чего его необходимо активировать нажатием левой кнопки.

\section*{Внутреннее устройство}

Внутреннее устройство мыши показано на рисунке \ref{fig:NumonicsManagerMouse2Inside}. Как можно видеть, она представляет собой оптомеханическую конструкцию (в отличие от механических энкодеров Manager Mouse первого поколения), но расположение энкодеров осталось прежним: диски оптических прерывателей закреплены непосредственно на осях колес. Физическое разрешение мыши увеличено в два раза по сравнению с моделью первого поколения и составляет 200 DPI \cite{resolution}.

\begin{figure}[h]
    \centering
    \includegraphics{1989_numonics_manager_mouse_cordless_2gen/inside_30.jpg}
    \caption{Manager Mouse в разобранном виде}
    \label{fig:NumonicsManagerMouse2Inside}
\end{figure}

Также на фото можно видеть блок из четырех перезаряжаемых никель-кадмиевых батарей, вставленных в съемный черный пластиковый контейнер, которые значительно увеличивают вес мыши. Для их зарядки и для обеспечения информационного обмена используется дополнительная печатная плата, расположенная поверх основной платы мыши в передней части корпуса: она содержит разъем для зарядки и два инфракрасных излучающих диода, расположенных под углом 60 градусов \cite{italianReview}.

\begin{thebibliography}{9}

\bibitem {torringtonAdd} The smallest link in your desktop publishing system is the strongest (advertising) // PC Magazine, V. 6, No. 17, October  13, 1987. - p. 305.\url{https://archive.org/details/PC-Mag-1987-10-13/page/304/mode/2up}

\bibitem {admaynard} This Little Fella Means Business [adv.] // PC Magazine, Vol. 5, No. 1, January 14, 1986. P. 37. \url{https://archive.org/details/PC-Mag-1986-01-14/page/n71/mode/2up}

\bibitem {adnumonics} Don't let your mouse get hung-up on itself [adv.] // PC Magazine, Vol. 7, No. 6, March 29, 1988. P. 282 // \url{https://archive.org/details/PC-Mag-1988-03-29/page/n277/mode/2up}

\bibitem {numonics} About Numonics Corporation -- numonics.com \url{https://web.archive.org/web/19981203122830/http://www.numonics.com/}

\bibitem {buxtonNumonicsCardboard} NumonicsCordlessMouse053 -- billbuxton.com \url{https://www.billbuxton.com/inputTimelineAssets/NumonicsCordlessMouse053.pdf}

\bibitem{italianReview} MC Microcomputer, No. 87, July-August 1989. P. 68. \url{https://archive.org/details/MC_microcomputer-087/page/n67/mode/2up}

\bibitem{resolution} Tout sur les souris // Compatibles PC Magazine, No. 29, Decembre 1989. P. 19 \url{https://archive.org/details/compatibles-pc-magazine-29-decembre-1989/page/n19/mode/2up}

%\bibitem {maynard} Hart G. Building a better mouse interface. Maynard mouse // PC Magazine, Vol. 5, No. 4, February 25, 1986. - p. 170--172 \url{https://archive.org/details/PC-Mag-1986-02-25/page/n177/mode/2up}

%\bibitem {numonicsOldwheels} Numonics corp. Manager mouse, Manager mouse cordless. //  PC Magazine, V. 7, No. 3, February 14, 1989. - p. 268. \url{https://archive.org/details/PC-Mag-1989-02-14/page/n267/mode/2up}

%\bibitem {buxtonTorrington} Manager Mouse Cordless \url{https://www.billbuxton.com/inputTimelineAssets/TorringtonBrochure.pdf}




\end{thebibliography}
\end{document}


